\documentclass[12pt]{article}
\usepackage[utf8]{inputenc}
\usepackage[T1]{fontenc}
\usepackage{hyperref}
\title{Eastern Tamang Language Resources for K-12 Curriculum Development}
\author{Compiled by Claude Sonnet 5 (claude-sonnet-5) for Andr\'as Kornai}
\date{}
\begin{document}
\maketitle
\section{Introduction}
Eastern Tamang (ISO 639-3 taj; Glottocode east2347) is a Sino-Tibetan language of Nepal with about 1,200,000 speakers. My searches found no resources specific to Eastern Tamang. The only materials found concern Western Tamang (a different ISO code). They are listed below as related material only, not as Eastern Tamang resources. They may help curriculum planners, but they cannot be assumed to match Eastern Tamang.
\section{Primary Educational Resources}
No Eastern Tamang primers, textbooks or mother-tongue curriculum documents were found in search results.
\section{Linguistic and Reference Materials}
\begin{itemize}
\item Western Tamang dictionary (Tamang--Nepali, with a Nepali--Tamang section; SIL International and Mother Tongue Center Nepal, 2017), web edition. This is Western Tamang, not Eastern: \url{https://www.webonary.org/wtamang/overview/?lang=en}
\item Same dictionary, download page: \url{https://www.webonary.org/wtamang/download/?lang=en}
\item Same dictionary, alphabet overview: \url{https://www.webonary.org/wtamang/overview/alphabet/}
\item Everitt, Taylor and Tamang (1972), a 40-page vocabulary that Glottolog files under Western Tamang: \url{https://glottolog.org/resource/reference/id/26744}
\item A Tribhuvan University sociolinguistic survey of Tamang in Gorkha district; I did not see its contents beyond the cover, and it is a general survey, not curriculum material: \url{https://elibrary.tucl.edu.np/bitstream/123456789/16320/1/Cover.pdf}
\end{itemize}
\section{Religious and Community Resources}
None found in search results.
\section{Digital and Audio Resources}
None found for Eastern Tamang. The Western Tamang dictionary is also offered as an app and PDF, according to the search summary.
\section{Curriculum Development Context}
The reference description says Eastern Tamang is taught as a subject. I found no source confirming this, and no Nepali government mother-tongue curriculum for Tamang appeared in the results. Further searching (Nepal Curriculum Development Centre mother-tongue materials, Tamang-language organisations) is needed before any curriculum work.
\end{document}